\section{Formazione relativa alle competenze trasversali}

``Frontiere dell'intelligenza artificiale'' di Giovanni Colavizza, 20 giugno 2025, 14:00-17:00.

``Artificial Intelligence Tools for Content Production in Academia'' di Paolo Torroni e Alessandro Gallo. Bologna, 26 maggio 2025, 23 giugno 2025, 30 giugno 2025, 15:00-18:00.

``Nuove macchine e antichi saperi: IA, Supercalcolo e dati per il patrimonio culturale''. Convegno su intelligenza artificiale, supercalcolo e dati per il patrimonio culturale. Frequentato online, 3 dicembre 2024, 9:00-13:00.

``Using Generative Artificial Intelligence for processing tweets'' di Elisavet Chantavaridou (Università di West Attica). Frequentato online, 6 dicembre 2024, 15:00-16:30.